\documentclass[12pt]{ctexbeamer}	%声明文档类型为beamer幻灯片
\mode <presentation>

%\usetheme{default}
%\usetheme{AnnArbor}
%\usetheme{Antibes}
%\usetheme{Bergen}
%\usetheme{Berkeley}
%\usetheme{Berlin}
%\usetheme{Boadilla}
%\usetheme{CambridgeUS}
\usetheme{Copenhagen}
%\usetheme{Darmstadt}
%\usetheme{Dresden}
%\usetheme{Frankfurt}
%\usetheme{Goettingen}
%\usetheme{Hannover}
%\usetheme{Ilmenau}
%\usetheme{JuanLesPins}
%\usetheme{Luebeck}
% \usetheme{Madrid}
% \usetheme{metropolis}
%\usetheme{Malmoe}
%\usetheme{Marburg}
%\usetheme{Montpellier}
%\usetheme{PaloAlto}
%\usetheme{Pittsburgh}
%\usetheme{Rochester}
%\usetheme{Singapore}
%\usetheme{Szeged}
%\usetheme{Warsaw}

%\usecolortheme{albatross}
\usecolortheme{beaver}
%\usecolortheme{beetle}
%\usecolortheme{crane}
%\usecolortheme{dolphin}
%\usecolortheme{dove}
%\usecolortheme{fly}
%\usecolortheme{lily}
%\usecolortheme{orchid}
%\usecolortheme{rose}
%\usecolortheme{seagull}
%\usecolortheme{seahorse}
%\usecolortheme{whale}
%\usecolortheme{wolverine}

%\beamersetaveragebackground{black!10}

\setbeamercovered{transparent}
\setbeamertemplate{navigation symbols}{}

\usefonttheme{professionalfonts}
\useinnertheme{circles}%{rectangles}
\setbeamertemplate{itemize item}{$\circledast$}%{\checkmark}


\setbeamertemplate{title page}
{
%  \vbox{}
  \vfill
  \begin{centering}
  \includegraphics[height=1.0cm]{figure/hundsun-logo.png}
        \vskip0.6em\par%
    \begin{beamercolorbox}[sep=8pt,center,shadow=true,rounded=true]{title}
      \usebeamerfont{title}\inserttitle\par%
      \ifx\insertsubtitle\@empty%
      \else%
        \vskip0.25em%
        {\usebeamerfont{subtitle}\usebeamercolor[fg]{subtitle}\insertsubtitle\par}%
      \fi%
    \end{beamercolorbox}%
    \vskip1em\par
    \begin{beamercolorbox}[sep=8pt,center]{author}
      \usebeamerfont{author}\insertauthor
    \end{beamercolorbox}
    \begin{beamercolorbox}[sep=8pt,center]{institute}
      \usebeamerfont{institute}\insertinstitute
    \end{beamercolorbox}
    \begin{beamercolorbox}[sep=8pt,center]{date}
      \usebeamerfont{date}\insertdate
    \end{beamercolorbox}\vskip0.5em
  \end{centering}
  \vfill
}

\hypersetup{pdfpagemode=FullScreen} % makes your presentation go automatically to full screen

\usepackage{fontspec,xunicode,xltxtra}

\usepackage{multimedia}	%让文档支持多媒体
\usepackage{graphics}	%让文档支持图片
\usepackage{hyperref}	%让文档支持超链接
\usepackage{booktabs}	%让文档支持三线表格
\usepackage{amsmath}	%ams可以让文档支持数学公式
\usepackage{amsfonts}
\usepackage{amssymb}
\usepackage{color}
\usepackage{graphicx,psfrag}
\usepackage{epsfig}

\title[期权知识分享]{期权基础知识及 OPlus 模块介绍}	%幻灯片标题
\author[Zengyu Zeng]{曾锃煜}

\institute[Hundsun]{资管，研发四部}	%作者单位
\date{\today}

\setcounter{tocdepth}{1}

\begin{document}
\begin{frame}
\titlepage
\end{frame}

\AtBeginSection[]{
  \frame{
    \frametitle{大纲}
    \tableofcontents[currentsection]
  }
}

\begin{frame}
\frametitle{大纲}
\tableofcontents
\end{frame}

%%%%%%%%%%%%%%%%%%%%%%%%%%%%%%%%%%%%%%%%%%%%%%%%%%%%%%%%%%%%%%%%%%%%%%%%%%%%%%%

\section{基本概念}

\begin{frame}{期权定义}
  \begin{itemize} %enumerate;description
    \item 约定价格（对赌）
    \item 权利金
    \item 保证金
    \item 行使权利
    \item 最后交易日、到期日、行权日、行权时间
    \item 交收日
  \end{itemize}
\end{frame}

\begin{frame}{期权定义}
  \begin{block}{时间节点}
    \begin{itemize}
      \item 到期月的第四个星期三
      \item 多出来的 30 分钟
    \end{itemize}
  \end{block}
  \begin{block}{违约}
    \begin{itemize}
      \item 证券交收
      \item 资金交收
    \end{itemize}
  \end{block}
\end{frame}

\begin{frame}{期权分类}
  \begin{block}{买卖方向}
    \begin{itemize}
      \item 认购期权
      \item 认沽期权
    \end{itemize}
  \end{block}
  \begin{block}{期权内涵价值}
    \begin{itemize}
      \item 实值期权
      \item 虚值期权
      \item 平值期权
    \end{itemize}
  \end{block}
\end{frame}

\begin{frame}{中国期权交易}
  \begin{block}{投资门槛}
    \begin{itemize}
      \item 托管与资产不少于 100 万 RMB
      \item 通过测试
      \item 有模拟经历
      \item 信用
    \end{itemize}
  \end{block}
\end{frame}

\begin{frame}{中国期权交易}
  \begin{block}{衍生品合约账户}
    \begin{itemize}
      \item 记录合约持仓
      \item 期权交易申报
      \item 行权申报
    \end{itemize}
    不可用于存放现货证券，行权的交割标的物通过普通证券账户完成。
  \end{block}
  \begin{block}{衍生品保证金账户}
    \begin{itemize}
      \item 权利金
      \item 卖方保证金
      \item 行权资金交收
    \end{itemize}
  \end{block}
\end{frame}

\begin{frame}{中国期权交易}
  \begin{block}{保证金计算}
    认购期权（short call） = 权利金前结算价 + Max\{21\% $\times$ 标的前收盘价 - 认购期权虚值，10\% $\times$ 标的前盘收盘价\} \\
    认沽期权（short put） = Min\{权利金前结算价 + Max\{19\% $\times$ 标的前收盘价 - 认沽期权虚值，10\% $\times$ 行权价\}，行权价\}\\
  \end{block}
  期权卖方需要在每个交易日末补足保证金，否则会被强制平仓。部分权利金也被放入保证金内。
\end{frame}

\begin{frame}{中国期权交易}
  国内目前有 ETF 期权与个股期权进行交易。此处以上证 50ETF 作为期权交易举例。标的即为 50ETF，Exchange Trading Fund。
  \begin{block}{上证 50 指数}
    \begin{enumerate}
      \item 50
      \item 加权
      \item 算法：$$\frac{a}{c} \times b$$a = 实时加权股票净值, b = 基准指数, c = 基准日加权股票净值
    \end{enumerate}
  \end{block}
\end{frame}

\begin{frame}{中国期权交易}
  \begin{block}{50ETF}
    \begin{itemize}
      \item 指数基金
      \item 上证 50 指数
      \item 交易成本
    \end{itemize}
  \end{block}
  \begin{block}{交易类型}
    有六个类型，分别是“买入”、“卖出”、“备兑”的开仓与平仓。
  \end{block}
\end{frame}

\begin{frame}{中国期权交易}
  \begin{block}{备兑交易}
    \begin{enumerate}
      \item T 日：标的券锁定；备兑开仓
      \item E 日：被行权；行权指派；锁定需要行权的合约标的，其余释放
      \item E + 1 日：准备交收
    \end{enumerate}
    \begin{itemize}
      \item <1-> 锁定证券
      \item <2-> 备兑开仓
      \item <3-> 解锁证券
      \item <4-> 备兑平仓
    \end{itemize}
  \end{block}
\end{frame}

\begin{frame}{中国期权交易}
  \begin{block}{卖出交易}
    \begin{enumerate}
      \item T 日：卖出开仓
      \item E 日：被行权；行权指派；收取需要被行权部分期权持仓保证金，其余保证金释放
      \item E + 1 日：准备交收，ETF 为 E + 2
    \end{enumerate}
  \end{block}
\end{frame}

\begin{frame}{中国期权交易}
  \begin{block}{买入交易}
    \begin{enumerate}
      \item T 日：标的券锁定；备兑开仓
      \item E 日：若不行权，则被平仓；若行权，则在 E + 1 日行权交收
      \item E + 1 日：行权交收，ETF 为 E + 2
    \end{enumerate}
  \end{block}
\end{frame}

\begin{frame}{中国期权交易}
  \begin{block}{双向持仓}
    \begin{itemize}
      \item 盘中
      \item 日终
    \end{itemize}
  \end{block}
\end{frame}

\begin{frame}{中国期权交易}
  \begin{block}{套利}
    \begin{enumerate}
      \item 隐含波动率
      \item 卖出高估期权
      \item 买入正常期权
    \end{enumerate}
  \end{block}
  \pause
  \begin{block}{隐含波动率}
    $$\frac{\partial V}{\partial t} + \frac{1}{2}\sigma^2 S^2\frac{\partial^2 V}{\partial S^2} + rS\frac{\partial V}{\partial S} - rV=0$$
    V：期权价格，t：时间，S：标的券价格，r：无风险利率，$\sigma$就是隐含波动率，解该方程即可。
  \end{block}
\end{frame}

%%%%%%%%%%%%%%%%%%%%%%%%%%%%%%%%%%%%%%%%%%%%%%%%%%%%%%%%%%%%%%%%%%%%%%%%%%%%%%%

\section{OPlus 系统期权模块}

\begin{frame}{开户及基础信息维护}
  包括期权信息维护、期权类别信息、行情信息的维护。
  \begin{columns}
    \begin{column}{0.60\textwidth}
      \epsfig{figure=figure/account-struct.PNG, width=\textwidth}
    \end{column}
    \begin{column}{0.40\textwidth}
      \begin{itemize}
        \item 公司维护
        \item 产品、单元维护
        \item 组和维护，期权与标的为 1：1
        \item 资金账号
        \item 股东维护，通过股东表查找资金账号
        \item 柜台同步
        \item 经纪商、柜台、营业部维护
        \item 费用维护
      \end{itemize}
    \end{column}
  \end{columns}
\end{frame}

\begin{frame}{开户及基础信息维护}
  \begin{block}{委托方向}
    \begin{itemize}
      \item 买入开平仓
      \item 卖出开平仓
      \item 行权认购、行权认沽
      \item 证券锁定、证券解锁
    \end{itemize}
  \end{block}
  \begin{block}{账号维护}
    股东账号不区分衍生品，资金账号区分衍生品。\\
    资金账号配置为个股期权，清算模式通过分仓通道区分。
  \end{block}
\end{frame}

\begin{frame}{开户及基础信息维护}
  \begin{columns}
    \begin{column}{0.70\textwidth}
      \begin{tabular}{ c | c | c | c }
        \hline
         & 经手费 & 结算费 & 过户费 \\
        \hline
        开平仓（含备兑） & $\sqrt{}$ & $\sqrt{}$ &  \\
        \hline
        行权及撤销 &  & $\sqrt{}$ & $\sqrt{}$ \\
        \hline
        锁定解锁 & & & \\
        \hline
      \end{tabular}
    \end{column}
    \begin{column}{0.30\textwidth}
      \begin{itemize}
        \item 日常结算费
        \item 行权结算费
        \item 期权经手费，股票 3 元，ETF 1.3 元
        \item 过户费（沪深不同，需特殊处理）
      \end{itemize}
    \end{column}
  \end{columns}
\end{frame}

\begin{frame}{风险控制}
  包括保证金监控, 限仓监控, 限购, 限制，归档等操作流程。大部分为复用其他模块的风控逻辑，仅新增“期权保证金控制”。
\end{frame}

\begin{frame}{交易管理}
  主要是期权交易、行权、标的券锁定、解锁功能。指令交易和策略交易暂不支持。
  \begin{block}{涉及的表}
    \begin{itemize}
      \item 期权委托表 jy\_toptentrust
      \item 期权成交表 jy\_toptrealdeal
      \item 现货组合持仓表 jy\_tunitstock，新增锁定标的字段
      \item 现货产品持仓表 jy\_tfundstock
      \item 股东信息表 yh\_tstockholder
      \item 资产信息表 jy\_tassetday
    \end{itemize}
  \end{block}
  OPlus 系统的“期权属性表”没有包括“隐含波动率”字段，是否可以考虑加入？
\end{frame}

\begin{frame}{交易管理}
  成本与盈亏将“认购”、“认沽”、“备兑”三者分开计算。
  \begin{block}{累计成本}
    \large{$cost = a \times (1 - \frac{sellAmount}{curAmount + sellAmount})$}\\
  \end{block}
  \begin{block}{累计收益}
    \large{$profit = oriProfit + sellPrice - sellCost - a \times \frac{sellAmount}{curAmount + sellAmount}$}
  \end{block}
  a = 期初成本 + 当日买入金额 + 买入费用；cost = 当前成本；profit = 当前累计收益；sellAmount = 卖出数量；curAmount = 当前数量；sellPrice = 卖出金额；sellCost = 卖出费用；oriProfit = 期初累计收益。
\end{frame}

\begin{frame}{交易管理}
  \begin{block}{浮动盈亏}
    \large{$obligor = -1 \times a - cost$\\}
    \large{$right = a - cost$}
  \end{block}
  a = 持仓数量 $\times$ 合约乘数 $\times$ 最新价；obligor：义务方；right：权利方；cost = 开仓成本。
\end{frame}

\begin{frame}{交易管理}
  \begin{block}{下委托}
    \begin{itemize}
      \item 支持个股期权开平仓、标的券锁定解锁，期权行权委托
      \item 开平仓及行权时前台传入期权组合， 锁定解锁传入标的组合
      \item 行权时，价格类型为限价，价格为行权价。锁定解锁价格类型不为必传，价格传0
    \end{itemize}
  \end{block}
  \begin{itemize}
    \item 按委托序号撤单
    \item 异常撤单
    \item 撤补单
  \end{itemize}
\end{frame}

\begin{frame}{清算}
  包括期权交易清算、行权清算、违约清算等内容。\\
\end{frame}

%%%%%%%%%%%%%%%%%%%%%%%%%%%%%%%%%%%%%%%%%%%%%%%%%%%%%%%%%%%%%%%%%%%%%%%%%%%%%%%

\begin{frame}
\centerline{\Large 谢谢!}
\end{frame}

\end{document}
